\honest{Probability is in the Mind}{Eliezer Yudkowsky}{2008}

The Mind Projection Fallacy, my subject yesterday, is E.~T. Jaynes's term, from his ``long and hard-fought battle against the accursèd frequentists.'' Jaynes held that probabilities express partial information: ignorance is a fact about a mind, not about the thing it is ignorant of. Here is the classic example, simplified.

You have a coin that you know is biased, but not which way or by how much. You flip it and cover it. Is the probability of heads 0.5? The frequentist says no: ``probability 0.5'' means the coin would come up heads half the time over infinitely many flips, and this coin is biased, so its probability of heads can be anything except 0.5. \nb{The speech merges two views: frequency over repeated trials (von Mises, who held that a single case has no probability) and a disposition of the whole set-up (Popper). On either, the claim is true.} The Bayesian says the coin has already landed, so a probability describes what I know, and theorems show that if I weigh my partial knowledge any other way I will make stupid bets; to plan, I weigh outcomes after heads no more than outcomes after tails. ``You can call that number whatever you like,'' but I will weigh heads and tails equally. \nb{As staged, the two do not disagree about any fact. One speaks of the coin's long-run frequency, the other of how to weigh this flip, and the Bayesian grants that the name is optional.} I side with the Bayesians.

Even a fair coin's 50 per cent ``may be just plain wrong,'' since a flip depends on how you hold and throw it and on the air; if you do not know which way it is biased this time, so what? \nb{Diaconis, Holmes and Montgomery (2007) built a machine that lands a coin heads ``one hundred percent of the time'' and found that natural flips come up as they started about 51 per cent of the time.} I believe a judge once asked, of a draft lottery said to be badly mixed, ``To whom is it unfair?'' \nb{Men did sue over the 1970 draft lottery, and the mixing was bad: men born late in the year drew more than their share of low numbers. The remark could not be traced.} A robot that watches an automatic flipper might predict it 90 per cent of the time. What is the real probability? There is none. The robot has one state of partial information and you another; the coin has no mind and assigns no probability, it just spins, bounces off air molecules and lands. \nb{This is the reference class problem, a known difficulty for frequentists. Their answer, which the post does not give, is that a probability is relative to a class of trials or a set-up, and the robot and you describe different ones.}

Now two brainteasers. A mathematician with two children, asked whether at least one is a boy, says yes: the prior chances of two boys, one of each and two girls are 1/4, 1/2 and 1/4; ``yes'' rules out two girls, leaving two boys at 1/3. Asked whether the eldest is a boy, ``yes'' gives 1/2, since the younger can be anything; the same for the youngest. But if at least one is a boy, it is the eldest or the youngest, so how can the first answer differ? With four cards, the aces and twos of hearts and spades, I draw two. ``At least one ace?'' ``Yes'' rules out the pair of twos among six equally likely pairs, so a pair of aces has 1/5. ``The ace of spades?'' ``Yes'' leaves three possible other cards, so 1/3; likewise for the ace of hearts. But the ace I hold must be one of the two, so how can the first answer be 1/5? All these calculations are correct. Different questions give different evidence, so you end in different states of knowledge. In Bayes's theorem, $P(H|E) = P(E|H)P(H)/P(E)$, dividing by $P(E)$ throws out what was eliminated and renormalizes. Before I answer, you expect ``yes'' to ``at least one ace'' with probability 5/6, but ``yes'' to ``the ace of spades'' with 1/2; for the children, 3/4 against 1/2. You are learning different things. The paradox comes only from thinking of probability as ``a property of cards,'' which would be like saying all green apples and all red apples weigh a pound but apples that are green or red weigh half a pound. \nb{A frequentist counting repeated draws gets the same 1/5 and 1/3. The view the apples defeat, that each set of cards carries its own probability as each apple has a weight, is not the frequentist's.}

``Probabilities express uncertainty, and it is only agents who can be uncertain.'' A blank map does not correspond to a blank territory. Ignorance is in the mind. \nb{Every example here is a classical system. The case that led Popper to propensities, quantum chance, is not discussed.}
